% GENERATED FILE. Edit canonical lessons, then run npm run generate:books.
% canonical-source-hash: fnv1a64:e7916c38eca2ac97
% canonical-lessons: IT-C329-rivolta, IT-C329-ribellarsi, IT-C329-invadere, IT-C329-arrendersi, IT-C329-liberare

\chapter{La rivolta}
\label{ch:it-la-rivolta}
\hlchaptermodality{329}

\begin{chapteropening}
\textbf{By the end of this chapter:} \emph{I can tell how people rose up: the revolt, rebelling, invading, surrendering and setting free.}

Complete the last lesson of chapter 329: I can tell how people rose up: the revolt, rebelling, invading, surrendering and setting free.
\end{chapteropening}

\section[la rivolta]{la rivolta — the revolt, the uprising}

\label{lesson:IT-C329-rivolta}



\begin{quote}
Before the new one: say the Italian for to govern, then the Italian for to found.
\end{quote}

\subsection*{You'll want to know: la rivolta}
\textbf{la rivolta} — \textquotedblleft{}the revolt, the uprising\textquotedblright{}.

\begin{quote}
La rivolta è cominciata in piazza.
\end{quote}

\begin{quote}
\emph{The uprising began in the square.}
\end{quote}

Built on \emph{rivoltare}, ‘to turn over’.

\begin{cousinweb}
\emph{rivolta} is built on \emph{rivoltare} \textquotedblleft{}to turn over\textquotedblright{} (EtymDB), from Latin \emph{volvere} \textquotedblleft{}to roll\textquotedblright{} (Lewis \& Short).
\end{cousinweb}

\begin{grammarlens}[title={What you've built}]
One more word for telling, explaining and describing at B1.
\end{grammarlens}

\subsection*{Your turn}
\begin{itemize}
\raggedright
  \item \emph{Say it:} \emph{la rivolta}
  \item \emph{Say it:} \emph{la rivolta}, once more
  \item \emph{Say it:} a sentence of your own with \emph{la rivolta}
  \item \emph{Recall:} say the Italian for to govern, then the Italian for to found, then say \emph{la rivolta} again
\end{itemize}

\begin{tcolorbox}[breakable,colback=teal!4,colframe=teal!35!black,title={Before you move on}]
What does la rivolta mean? (\textquotedblleft{}The revolt, the uprising\textquotedblright{}.) Say it once more.
\end{tcolorbox}
\section[ribellarsi]{ribellarsi — to rebel}

\label{lesson:IT-C329-ribellarsi}



\begin{quote}
Before the new one: say the Italian for to found, then the Italian for the revolt.
\end{quote}

\subsection*{You'll want to know: ribellarsi}
\textbf{ribellarsi} — \textquotedblleft{}to rebel\textquotedblright{}.

\begin{quote}
I contadini si sono ribellati.
\end{quote}

\begin{quote}
\emph{The peasants rebelled.}
\end{quote}

\emph{Ribellarsi a} qualcuno.

\begin{cousinweb}
\emph{ribellarsi} is built on \emph{ribelle} \textquotedblleft{}rebel\textquotedblright{}, from Latin \emph{rebellis}, from \emph{rebellāre} \textquotedblleft{}to make war again\textquotedblright{}, built on \emph{bellum} \textquotedblleft{}war\textquotedblright{} (Lewis \& Short).
\end{cousinweb}

\begin{grammarlens}[title={What you've built}]
One more word for telling, explaining and describing at B1.
\end{grammarlens}

\subsection*{Your turn}
\begin{itemize}
\raggedright
  \item \emph{Say it:} \emph{ribellarsi}
  \item \emph{Say it:} \emph{ribellarsi}, once more
  \item \emph{Say it:} a sentence of your own with \emph{ribellarsi}
  \item \emph{Recall:} say the Italian for to found, then the Italian for the revolt, then say \emph{ribellarsi} again
\end{itemize}

\begin{tcolorbox}[breakable,colback=teal!4,colframe=teal!35!black,title={Before you move on}]
What does ribellarsi mean? (\textquotedblleft{}To rebel\textquotedblright{}.) Say it once more.
\end{tcolorbox}
\section[invadere]{invadere — to invade}

\label{lesson:IT-C329-invadere}



\begin{quote}
Before the new one: say the Italian for the revolt, then the Italian for to rebel.
\end{quote}

\subsection*{You'll want to know: invadere}
\textbf{invadere} — \textquotedblleft{}to invade\textquotedblright{}.

\begin{quote}
L'esercito ha invaso il paese vicino.
\end{quote}

\begin{quote}
\emph{The army invaded the neighbouring country.}
\end{quote}

Past participle \emph{invaso}.

\begin{cousinweb}
\emph{invadere} is inherited from Latin \emph{invādere} \textquotedblleft{}to go into, to attack\textquotedblright{} (EtymDB).
\end{cousinweb}

\begin{grammarlens}[title={What you've built}]
One more word for telling, explaining and describing at B1.
\end{grammarlens}

\subsection*{Your turn}
\begin{itemize}
\raggedright
  \item \emph{Say it:} \emph{invadere}
  \item \emph{Say it:} \emph{invadere}, once more
  \item \emph{Say it:} a sentence of your own with \emph{invadere}
  \item \emph{Recall:} say the Italian for the revolt, then the Italian for to rebel, then say \emph{invadere} again
\end{itemize}

\begin{tcolorbox}[breakable,colback=teal!4,colframe=teal!35!black,title={Before you move on}]
What does invadere mean? (\textquotedblleft{}To invade\textquotedblright{}.) Say it once more.
\end{tcolorbox}
\section[arrendersi]{arrendersi — to surrender, to give up}

\label{lesson:IT-C329-arrendersi}



\begin{quote}
Before the new one: say the Italian for to rebel, then the Italian for to invade.
\end{quote}

\subsection*{You'll want to know: arrendersi}
\textbf{arrendersi} — \textquotedblleft{}to surrender, to give up\textquotedblright{}.

\begin{quote}
Dopo un mese la città si è arresa.
\end{quote}

\begin{quote}
\emph{After a month the city surrendered.}
\end{quote}

Past participle \emph{arreso}. \emph{Non arrenderti!} — do not give up!

\begin{cousinweb}
\emph{arrendersi} is \emph{a-} with \emph{rendere} \textquotedblleft{}to give back\textquotedblright{}, from Vulgar Latin \emph{rendere}, a form of Latin \emph{reddere} \textquotedblleft{}to give back\textquotedblright{} (Lewis \& Short).
\end{cousinweb}

\begin{grammarlens}[title={What you've built}]
One more word for telling, explaining and describing at B1.
\end{grammarlens}

\subsection*{Your turn}
\begin{itemize}
\raggedright
  \item \emph{Say it:} \emph{arrendersi}
  \item \emph{Say it:} \emph{arrendersi}, once more
  \item \emph{Say it:} a sentence of your own with \emph{arrendersi}
  \item \emph{Recall:} say the Italian for to rebel, then the Italian for to invade, then say \emph{arrendersi} again
\end{itemize}

\begin{tcolorbox}[breakable,colback=teal!4,colframe=teal!35!black,title={Before you move on}]
What does arrendersi mean? (\textquotedblleft{}To surrender, to give up\textquotedblright{}.) Say it once more.
\end{tcolorbox}
\section[liberare]{liberare — to free, to set free}

\label{lesson:IT-C329-liberare}



\begin{quote}
Before the new one: say the Italian for to invade, then the Italian for to surrender.
\end{quote}

\subsection*{You'll want to know: liberare}
\textbf{liberare} — \textquotedblleft{}to free, to set free\textquotedblright{}.

\begin{quote}
Alla fine hanno liberato i prigionieri.
\end{quote}

\begin{quote}
\emph{In the end they freed the prisoners.}
\end{quote}

Also to clear: \emph{liberare la stanza}.

\begin{cousinweb}
\emph{liberare} is inherited from Latin \emph{līberāre} \textquotedblleft{}to set free\textquotedblright{} (EtymDB), from \emph{līber} \textquotedblleft{}free\textquotedblright{}.
\end{cousinweb}

\begin{grammarlens}[title={What you've built}]
One more word for telling, explaining and describing at B1.
\end{grammarlens}

\subsection*{Your turn}
\begin{itemize}
\raggedright
  \item \emph{Say it:} \emph{liberare}
  \item \emph{Say it:} \emph{liberare}, once more
  \item \emph{Say it:} a sentence of your own with \emph{liberare}
  \item \emph{Recall:} say the Italian for to invade, then the Italian for to surrender, then say \emph{liberare} again
\end{itemize}

\begin{tcolorbox}[breakable,colback=teal!4,colframe=teal!35!black,title={Before you move on}]
What does liberare mean? (\textquotedblleft{}To free, to set free\textquotedblright{}.) Say it once more.
\end{tcolorbox}
